% Figure: The omics stack: genome, transcriptome, proteome, metabolome,
% lipidome, fluxome, showing what each layer measures and at what
% scale.
\begin{figure}[htbp]
\centering
\begin{tikzpicture}[
  layer/.style={draw=engblack, rounded corners=2pt, minimum width=4.0cm,
    minimum height=0.9cm, align=center, font=\small},
  ann/.style={anchor=west, font=\footnotesize, color=enggray},
  arrow/.style={->, thick, draw=engblack!60},
]

% Vertical stack of layers (top: genome; bottom: phenotype)
\node[layer, fill=engblue!10]   (g) at (0, 5.0) {Genome};
\node[layer, fill=engblue!20]   (t) at (0, 3.7) {Transcriptome};
\node[layer, fill=englime!20]   (p) at (0, 2.4) {Proteome};
\node[layer, fill=engamber!20]  (m) at (0, 1.1) {Metabolome / Lipidome};
\node[layer, fill=engred!15]    (f) at (0, -0.2) {Fluxome};
\node[layer, fill=enggray!20]   (ph) at (0, -1.5) {Phenotype};

% Right-side annotations: what is measured and how
\node[ann, align=left] at (3.0, 5.0)
  {$\sim 3 \times 10^{9}$~bp; static.\\
   Measured by short-read sequencing,\\
   stored as FASTA + VCF.};
\node[ann, align=left] at (3.0, 3.7)
  {$\sim 2 \times 10^{4}$~genes, $10^{5}$~transcript isoforms.\\
   Dynamic, condition-specific.\\
   Measured by RNA-seq, scRNA-seq.};
\node[ann, align=left] at (3.0, 2.4)
  {$\sim 2 \times 10^{4}$~proteins, $10^{6}$~proteoforms.\\
   Post-translational state.\\
   Measured by mass spectrometry.};
\node[ann, align=left] at (3.0, 1.1)
  {$\sim 10^{4}$~metabolites, $10^{3}$~lipid species.\\
   Dynamic, fast-turnover.\\
   Measured by LC-MS, GC-MS, NMR.};
\node[ann, align=left] at (3.0, -0.2)
  {$\sim 10^{3}$~reactions in central metabolism.\\
   Inferred, not measured directly.\\
   FBA, $^{13}$C tracing, kinetic models.};
\node[ann, align=left] at (3.0, -1.5)
  {Growth, survival, disease, fitness.\\
   The reason the rest of the stack exists.};

% Vertical arrows
\draw[arrow] (g.south) -- (t.north);
\draw[arrow] (t.south) -- (p.north);
\draw[arrow] (p.south) -- (m.north);
\draw[arrow] (m.south) -- (f.north);
\draw[arrow] (f.south) -- (ph.north);

% Left-side process labels
\node[font=\footnotesize\itshape, color=enggray, anchor=east, align=right]
  at (-2.2, 4.35) {transcription\\+ regulation};
\node[font=\footnotesize\itshape, color=enggray, anchor=east, align=right]
  at (-2.2, 3.05) {translation\\+ PTM};
\node[font=\footnotesize\itshape, color=enggray, anchor=east, align=right]
  at (-2.2, 1.75) {enzymatic\\activity};
\node[font=\footnotesize\itshape, color=enggray, anchor=east, align=right]
  at (-2.2, 0.45) {flux distribution\\through network};
\node[font=\footnotesize\itshape, color=enggray, anchor=east, align=right]
  at (-2.2, -0.85) {organism-level\\integration};

\end{tikzpicture}
\caption[The omics stack]{The omics stack reads bottom-up as the
chain of causation by which genome becomes phenotype, and top-down
as the chain of measurement by which engineers reconstruct what
the cell is doing. Each layer has its own measurement technology,
its own scale, its own dynamic range, and its own characteristic
artefact class. A study that measures only one layer measures the
projection of a many-dimensional state onto a single axis; the
multi-omics integration problem is the engineering problem of
estimating the joint state from partial views with very different
noise statistics. Counts and figures are order-of-magnitude human
estimates current as of 2024.}
\label{fig:vol06:ch11:omics-stack}
\end{figure}
